\section{Control of physical network competitors}
\label{app:physical}

The sharp code-atom theorem and its quantitative support bound control
all physical network competitors, including those whose outputs have
Fourier components outside the logical code. We now prove that every
minimum-cost feasible score is the cheapest live rectangle.

Here are the physical definitions, allowing the column code of
Section~\ref{sec:cycles}. Take $\ell$ disjoint rail pairs $B_j^0,B_j^1$
and nonzero columns $a_j\in\mathbb F_2^m$ spanning the logical space,
together with an independent shared carrier $x_\star$.
For a uniformly random physical input put
\[
 y_j^a=\chi_{B_j^a}(x),\qquad b=x_\star\prod_j y_j^0,\qquad
 (-1)^{\zeta_j}=y_j^0y_j^1,\qquad z=\sum_j a_j\zeta_j.
\]
The logical vector $z$ is uniform on $G=\mathbb F_2^m$.
The code space $\mathcal C$ is the span of physical characters
$\chi_{A_s}=b(-1)^{s\cdot z}$, where
$A_s=\{\star\}\cup\bigcup_jB_j^{a_j\cdot s}$. The odd downward
Fourier space $\mathcal D$ is spanned by odd-degree characters with
support contained in some $A_s$. All physical rail sizes are even. The independent-rail code is the special case $a_j=e_j$.
All projections and norms are in the uniform physical $L^2$ space.
In particular,
\[
 P_{\mathcal C}F=b\,g(z),\qquad g(z)=\mathbb E[bF\mid z].
\]
The target is $f(x)=b(x)$ when $z=0$, and $f(x)=-b(x)$ otherwise.
Each revealed logical fiber means that all its physical inputs have
been labeled. Closed global margins mean $|F(x)|\ge1/2$ everywhere;
correct labels on revealed fibers mean $f(x)F(x)\ge1/2$ there.


Let $M=2^m$, $m\ge6$. Write
$w_s\ge w_0>0$ for code weights and
$W(z)=M^{-1}\sum_s w_s(-1)^{s\cdot z}$. Set $u=e_1+e_2$.
Take an affine coset $U=v_0+H$ with $K_U=|U|\le M/32$,
where $u\notin H$ and $v_0\notin H+\langle u\rangle$. Put
$H'=H+\langle u\rangle$, $U'=U+\langle u\rangle$, and
$E'=H'^\perp\setminus(H'+\langle v_0\rangle)^\perp$.
Assume $U'\cup\{u\}$ is independent in the Cayley graph whose
edges add one of the columns $a_j$. In the independent-rail case,
one can use the earlier coset with its first $q\ge4$ bits fixed and
odd total parity. The cycle coset also satisfies these conditions.
Now $|E'|=M/(4K_U)\ge8$, every $s\in E'$ has $s\cdot u=0$, and
\[
 \sum_{s\in E'}(-1)^{s\cdot z}
 =|E'|(\mathbf1_{H'}(z)-\mathbf1_{U'}(z)).
\]

For every nonzero coset $c$ of $\langle u\rangle$, put
\[
 k_c=\max\{0,W(z),W(z+u)\}\quad(z\in c),\qquad
 K_* =\sum_{c\ne0}k_c,\qquad w_*=W(u).
\]
Assume
\begin{equation}\label{eq:rect-dominance}
 0<w_*\le w_0/16,\qquad K_*\le w_*/4,\qquad M w_*\ge16w_0,
 \qquad W(v+u)>W(v)\ge0\quad(v\in U).
\end{equation}

\subsection{The common dual}

Put $L=M/2$ and define
\[
 h=\frac{(L-1)(w_*-4w_0/M)+K_*}{L-2},\qquad
 b_* =\frac{M(h+K_*)}{L-1},\qquad a_*=2w_0-b_*.
\]
These satisfy $3w_*/4\le h\le5w_*/4$, $b_*\le4w_*\le w_0/4$,
and $7w_0/4\le a_*\le2w_0$. Set
\[
 r_c=\frac{(L-1)k_c}{Lh+K_*},\qquad R_*=\sum_{c\ne0}r_c\le1/3.
\]
On $u^\perp$ define the strictly positive probability distribution
\[
 P(s)=\frac1L\left(1+\sum_{c\ne0}r_c(-1)^{s\cdot z_c}\right).
\]
Its mass at zero is
\[
 p=P(0)=\frac{h+K_*}{Lh+K_*},\qquad
 1-p=\frac{(L-1)h}{Lh+K_*}.
\]
Let $Q$ be $P$ conditioned on $s\ne0$, extended by zero outside
$u^\perp\setminus\{0\}$. Define a logical dual by
\begin{equation}\label{eq:rect-dual}
 \widehat\lambda(s)=w_s-2w_0\mathbf1_{s=0}
                    -MhQ(s)-\frac{a_*}{|E'|}\mathbf1_{s\in E'}.
\end{equation}
For every $s$, $|\widehat\lambda(s)|\le w_s$. Indeed its value at zero
is $-w_0$, and it equals $w_s$ if $s\cdot u=1$. Otherwise
$MhQ(s)\le Mh p/(1-p)=b_*$, so the total subtraction is at most
$b_*+a_*/|E'|\le2w_0\le2w_s$.

For $z\notin\{0,u\}$, Fourier orthogonality gives
\[
 \mathbb E_Q(-1)^{s\cdot z}
 =\frac{k_{[z]}}h-\frac{h+K_*}{(L-1)h}.
\]
Consequently
\begin{equation}\label{eq:rect-dual-values}
 \lambda(z)=M(W(z)-k_{[z]})-a_*
            -a_*\mathbf1_{H'}(z)+a_*\mathbf1_{U'}(z)
 \quad(z\ne0,u).
\end{equation}
The defining equation for $h$ also gives
\[
 \lambda(u)=0,\qquad \lambda(0)=M(W(0)-w_*)\ge a_*.
\]
In particular,
\[
 \lambda(v)=M\nu_v<0,\qquad \lambda(v+u)=0,
 \qquad \nu_v:=W(v)-W(v+u)\quad(v\in U).
\]
On the observed logical region
$\mathcal O=G\setminus(U'\cup\{u\})$, its sign is positive at zero
and negative elsewhere, with magnitude at least $a_*$.

There is also a useful strict coefficient slack. If $s\ne0$ and
$s\cdot u=0$, then
\begin{equation}\label{eq:rect-invariant-slack}
 \min\{w_s-\widehat\lambda(s),w_s+\widehat\lambda(s)\}\ge w_*.
\end{equation}
For the first term use
$MhQ(s)\ge2h(1-R_*)/(1-p)\ge w_*$.
For the second use
$2w_s-b_*-a_*/|E'|\ge a_*(1-1/|E'|)>w_*$.

\subsection{Observability outside a logical independent set}

\begin{lemma}\label{lem:independent-observability}
If $J\subseteq G$ is independent in the Cayley graph generated by
the columns $a_j$,
then every $A\in\mathcal D\cap\mathcal C^\perp$ satisfies
\[
 \|A\|_2^2\le2\mathbb E[\mathbf1_{z\notin J}A^2].
\]
\end{lemma}
\begin{proof}
Every Fourier support $T$ contributing to $A$ is an odd-degree proper
subset of a carrier-containing transversal. It cannot occupy a full
rail at every pair: with the carrier it would be a retained code
character, while without the carrier it would have even degree.
Assign it to the first rail pair $i$ at which its restriction is neither
complete rail. Write $A=\sum_i A_i$ for this partition.

Conditioned on a logical fiber, two physical characters have nonzero
inner product only when their symmetric difference is a union of full
rail pairs with indicator vector $(a_j\cdot s)_j$ for some $s$.
A symmetric difference containing the carrier has zero conditional
expectation, since the carrier is uniform and independent on each
logical fiber. If one support is a proper subset of a rail at pair $i$,
its complement in that pair cannot be contained in a single rail.
Thus conditional covariance with another downward character requires
the same restriction at pair $i$. It follows that $A_i,A_j$ are
conditionally orthogonal whenever $i\ne j$, and that
$\ell_i(z):=\mathbb E[A_i^2\mid z]$ is invariant under $z\mapsto z+a_i$.
Indeed every surviving logical Fourier character has $a_i\cdot s=0$.

An independent set contains at most one endpoint of every edge that
adds $a_i$.
Therefore $\mathbb E[\mathbf1_{z\in J}\ell_i(z)]\le\|A_i\|_2^2/2$.
Sum over $i$, using conditional orthogonality, to obtain the assertion.
\end{proof}

By the assumed Cayley independent-set condition,
Lemma~\ref{lem:independent-observability} applies to $\mathcal O$.

\subsection{Physical variance control}

Suppose a score $F$ and a nonnegative cost $B\le B_{\rm cap}$ satisfy
\[
 Q(F):=\sum_s w_s|\widehat F(A_s)|\le B,\qquad
 \|F\|_\infty\le H_{\rm cap},\qquad
 \|F-P_{\mathcal D}F\|_2^2\le S_{\rm at} B D,
\]
where $D=B-\langle b\lambda,F\rangle\ge0$ and $b$ is the physical
base character. These bounds follow from the sharp atom theorem and its support
stability estimate, namely
\[
 Q(\phi_a)\le\|a\|_2,\qquad
 \|\phi_a-P_{\mathcal D}\phi_a\|_2^2
 \le S_{\rm at}\|a\|_2(\|a\|_2-Q(\phi_a)),
 \quad \phi_a(x)=\tanh(a\cdot\bar x),
\]
for a representation $F=\sum_jc_j\phi_{a_j}$ of cost
$B=\sum_j|c_j|\|a_j\|_2$. Indeed the dual inequality gives
$|\langle b\lambda,\phi_a\rangle|\le Q(\phi_a)$, and
Cauchy--Schwarz bounds the squared sum of residuals by $S_{\rm at}BD$.
The elementary bound $|\phi_a(x)|\le\sqrt{d+1}\|a\|_2$ also gives
$H_{\rm cap}=\sqrt{d+1}B_{\rm cap}$. The atom theorem supplies
a polynomial value of $S_{\rm at}$.
Assume closed full-cube margins and correct labels outside $U$, with
logical target positive only at zero. Extra fully revealed fibers in
$U$ have the negative label. Write $P_{\mathcal C}F=bg$ and
\[
 S=\mathbb E[\mathbf1_{z\in\mathcal O}|\lambda(z)|
                         (fF-1/2)]\ge0.
\]
On each observed fiber, the variable $Y=fF$ is at least $1/2$.
Thus $\mathbb E[|Y-\mathbb EY|\mid z]\le
2(\mathbb E[Y\mid z]-1/2)$. Since $|\lambda|\ge a_*$ there,
this proves
$\mathbb E[\mathbf1_{\mathcal O}|F-P_{\mathcal C}F|]\le2S/a_*$.
Applying independent-set observability to
$P_{\mathcal D}F-P_{\mathcal C}F$ therefore yields
\begin{equation}\label{eq:rect-variance}
 L_F:=\|F-P_{\mathcal C}F\|_2^2
 \le C_{\rm phys}(D+S),\qquad
 C_{\rm phys}=\max\{16H_{\rm cap}/a_*,\ 5S_{\rm at}B_{\rm cap}\}.
\end{equation}
To see the constants, write $R=F-P_{\mathcal C}F=A+E$, where
$A=P_{\mathcal D}F-P_{\mathcal C}F$ and $E=F-P_{\mathcal D}F$.
The preceding conditional bound and $|R|\le2H_{\rm cap}$ give
$\mathbb E[\mathbf1_{\mathcal O}R^2]\le4H_{\rm cap}S/a_*$.
Observability and $(R-E)^2\le2R^2+2E^2$ imply
$\|A\|_2^2\le4\mathbb E[\mathbf1_{\mathcal O}R^2]+4\|E\|_2^2$.
Adding $\|E\|_2^2$ proves \eqref{eq:rect-variance}.

\subsection{Linear rounding to zero or a live rectangle}

Put $p_v=g(v)+1/2$ for $v\in U$ and
$\delta=D/w_*+S/a_*$. By \eqref{eq:rect-invariant-slack}, the nonconstant
$u$-invariant part $r$ of $g$ satisfies $\|r\|_\infty\le D/w_*$.
The region $\mathcal O_0=G\setminus(U'\cup\{0,u\})$ is invariant
under $u$, has density at least $1/2$, and consists of known negative
fibers. Averaging $g+1/2$ there eliminates the $u$-anti-invariant part.
Consequently
\[
 |\widehat g(0)+1/2|\le D/w_*+2S/a_*,\qquad
 \sup_z|g(z)+g(z+u)+1|\le4\delta.
\]
Each partner $v+u$ is a known negative fiber, so $p_v\ge-4\delta$.
In particular the total negative mass of $p$ is at most $M\delta$.

Compare $g$ with
\[
 g_{\rm ref}=-1/2+\mathbf1_0-\mathbf1_u
              +\sum_{v\in U}p_v(\mathbf1_v-\mathbf1_{v+u}).
\]
The logical $\ell^1$ norm of their difference is at most
\[
 2MS/a_*+4(K_U+1)\delta\le3M\delta.
\]
Here the observed fibers contribute at most $MS/a_*$; the fiber $u$
contributes at most the anchor excess plus $4\delta$; and the $K_U$
partner fibers contribute at most $4K_U\delta$.

There are characters $s$ with $s\cdot u=1$ and
$(-1)^{s\cdot v}=-1$ for every $v\in U$: prescribe the two independent linear conditions
$s\cdot u=1$ and $s\cdot v_0=1$ within $H^\perp$. For any such $s$, the dual is saturated on its
positive side, so $\widehat g(s)\ge-D/(2w_s)$. Since
\[
 \widehat g_{\rm ref}(s)=\frac2M\left(1-\sum_{v\in U}p_v\right),
\]
we obtain
\[
 \sum_{v\in U}p_v\le1+\frac{MD}{4w_0}+\frac32M\delta
 \le1+2M\delta.
\]
Thus the total positive mass is at most $1+3M\delta$.

On a fiber with $0<p_v<1$, the mean $g(v)=p_v-1/2$ lies
between the two allowed output intervals. A variable $Y=bF$ with
$|Y|\ge1/2$ and mean $g(v)$ has variance at least
$1/4-g(v)^2=p_v(1-p_v)$. Consequently
\[
 V_{\rm frac}:=\sum_{0<p_v<1}p_v(1-p_v)\le M L_F.
\]
Round each $p_v$ to its nearest $b_v\in\{0,1\}$. Negative coordinates
contribute at most $M\delta$ to the rounding distance. Coordinates above
one contribute at most $3M\delta$, because total positive mass is at most
$1+3M\delta$. The fractional coordinates contribute at most $2V_{\rm frac}$.
Hence $A:=\|p-b\|_1\le4M\delta+2V_{\rm frac}$.
If $b=0$, choose the zero vector. Otherwise choose one coordinate with
$b_v=1$; it is live because removed fibers have $p_v\le0$. The distance
from $b$ to this point mass is at most $3M\delta+A$. Therefore
\begin{equation}\label{eq:rect-rounding}
 \operatorname{dist}_1(p,\{0\}\cup\{e_v:v\text{ live}\})
 \le11M\delta+4V_{\rm frac}
 \le B_{\rm round}(D+S),
\end{equation}
where
\[
 B_{\rm round}=M(11/w_*+11/a_*+4C_{\rm phys}).
\]

\subsection{Absorption and classification of every minimum}

The baseline score $g_0=-1/2+\mathbf1_0-\mathbf1_u$ and rectangle scores
$g_v=g_0+\mathbf1_v-\mathbf1_{v+u}$ saturate the common dual. Their
weighted Fourier costs are
\[
 C_0=w_0/2+W(0)-W(u),\qquad C_v=C_0+\nu_v.
\]
The exact calibrated decomposition is
\[
 B=C_0+D+S+\sum_{v\in U}\nu_v p_v.
\]
No omitted-fiber term remains: $\lambda(u)=0$ and
$\lambda(v+u)=0$ for all $v\in U$.
Let $W_{\rm eff}=\max_{v\in U}|\nu_v|$. If
$W_{\rm eff}B_{\rm round}<1$, rounding gives
\[
 B\ge C_{\rm choice}+(1-W_{\rm eff}B_{\rm round})(D+S),
\]
where the choice is zero or a live point mass. The zero choice has cost
$C_0$, strictly greater than every live star because $\nu_v<0$.
Consequently any minimum-cost completion is a cheapest live rectangle.
If its effective costs are pairwise distinct, the score is unique:
equality forces $D=S=0$, then $L_F=0$, exact anti-invariance around
$-1/2$, and the rounded point mass determine every logical and physical
value.

Lemma~\ref{lem:dominant-increments} and the cycle construction make
\eqref{eq:rect-dominance} hold for large $m$. The atom theorem gives
$S_{\rm at}=2\max(2,96\sqrt6\,W^2)=O(n_0^2)$, and one may take
$B_{\rm cap}=O(n_0)$. Every factor in $B_{\rm round}$ except $2^m$
is therefore polynomial in $m$.
The selected prefix factors give the required small Fourier costs. On $(0,t_0]$, the Laplace kernel is bounded and
$A_{\rm prefix}(t)\le(2t)^P$, while the other coefficient factors
are at most one. Therefore
\[
 |\nu_v|\le O(2^P P!/n_0^{P+1})+\exp(-\Omega(n_0)),
\]
uniformly over its affine live envelope. Since $P=\Theta(m)$ and
$n_0=\Theta(m^3)$, this again absorbs $2^m$ times any polynomial.
The explicit polynomial atom stability is what permits this
polynomial-resource choice.

This proves the calibration statement for the cycle construction.
Section~\ref{sec:cycles} supplies its distinct effective costs, and
Appendix~\ref{app:posterior} transfers the score-space cost gaps to
the actual fixed-width Gaussian posterior.
